% Generated draft: requires mathematical and source review.
\documentclass[11pt]{article}
\usepackage[margin=25mm]{geometry}
\usepackage{fontspec}
\setmainfont{Latin Modern Roman}
\usepackage{amsmath,amssymb,mathtools,unicode-math}
\setmathfont{Latin Modern Math}
\usepackage{enumitem,needspace,xurl,hyperref}
\setlength{\parindent}{0pt}
\setlength{\parskip}{5pt}
\setlength{\emergencystretch}{3em}
\newcommand{\N}{\mathbb{N}}
\newcommand{\Z}{\mathbb{Z}}
\newcommand{\Q}{\mathbb{Q}}
\newcommand{\R}{\mathbb{R}}
\newcommand{\C}{\mathbb{C}}
\newcommand{\F}{\mathbb{F}}
\newcommand{\A}{\mathbb{A}}
\newcommand{\PP}{\mathbb{P}}
\newcommand{\E}{\mathbb{E}}
\newcommand{\eps}{\varepsilon}
\newcommand{\dd}{\,\mathrm{d}}
\begin{document}
\section*{UnsolvedMath ID: 30003318}
\textbf{Source problem code:} catalogue entry
\subsection{Definitions and mathematical statement}
An ordered field \(K\) is real closed if every positive element has a square root in \(K\) and every polynomial in \(K[X]\) of odd degree has a root in \(K\). An integer part of \(K\) is a subring \(I\subseteq K\), containing the same identity element, whose induced order is discrete with least positive element \(1\), and such that
\[
\forall x\in K\;\exists a\in I\quad a\leq x<a+1.
\]
Write \(I_{\geq 0}=\{a\in I:a\geq 0\}\). In the arithmetic language \(L_{\mathrm{ar}}=\{0,1,+,\cdot,<\}\), true arithmetic is the complete theory
\[
\operatorname{Th}_{L_{\mathrm{ar}}}(\N,0,1,+,\cdot,<).
\]
The phrase ``an integer part that is a model of true arithmetic'' is interpreted here in the usual arithmetic sense that its nonnegative cone, with the inherited operations and order, satisfies this theory. This convention requires confirmation against the intended source formulation.

The supplied record groups three questions:
\begin{enumerate}
\item Investigate the behavior of fast-growing functions, including the Ackermann function, in real closed fields. The record does not specify a version of the Ackermann function, a domain for such functions, a method of extending them beyond \(\N\), or a property to prove. This part therefore remains a research direction rather than a determinate mathematical assertion.
\item Let \(\mathbf{No}\) denote the ordered field of surreal numbers. Does there exist an integer part \(I\) of \(\mathbf{No}\) such that
\[
(I_{\geq 0},0,1,+,\cdot,<)
\models
\operatorname{Th}_{L_{\mathrm{ar}}}(\N,0,1,+,\cdot,<)?
\]
Here the integer-part condition is the condition above with \(K=\mathbf{No}\). Since \(\mathbf{No}\) is a proper class, the permitted class framework and the meaning of satisfaction for the resulting class structure must be specified before this can be treated as an ordinary formal existence statement. The supplied record does not specify that framework.
\item Does every real closed field \(K\) have a normal integer part? The universal quantifier ranges over all real closed fields, with no additional hypothesis supplied. The adjective ``normal'' is not defined in the record. Consequently, this question cannot be made into a fully specified predicate on integer parts from the supplied evidence; its intended additional condition must be recovered from the source.
\end{enumerate}
These are statement-only formulations with unresolved qualifications, not mathematical resolutions.
\subsection{Short English statement}
The problem asks how fast-growing arithmetic functions relate to real closed fields, whether the surreal numbers have an integer part whose nonnegative cone satisfies true arithmetic, and whether every real closed field has a normal integer part. The fast-growth target, the meaning of ``normal'', and the class framework for the surreal question remain unspecified.
\subsection{Sources}
\begin{itemize}
\item Supplied local UnsolvedMath registry record, permanent ID \(30003318\), title ``Arithmetic Integer Parts of Surreal and Real Closed Fields'', source problem code \texttt{catalogue entry}; project registry locations: \texttt{data/unsolvedmath/problems.json} and \texttt{catalog/problems.json}. No individual proposer is identified in the supplied record.
\item Supplied report citation: \emph{Mini-Workshop: Surreal Numbers, Surreal Analysis, Hahn Fields and Derivations} (2017), \emph{Oberwolfach Reports}, pp.\ 3313--3372, DOI \texttt{10.4171/owr/2016/60}. \url{https://doi.org/10.4171/owr/2016/60}. The registry separately gives proposed year 2016 and the tag \texttt{owr-year:2016}; these supplied dates are retained without reconciling their roles.
\item Supplied literature reference, labelled ``Carl--Krapp'': \url{https://doi.org/10.4115/jla.2021.13.3}.
\item Supplied literature reference, labelled ``Krapp thesis'': \url{https://d-nb.info/1202012558/34}.
\item Supplied literature reference, labelled ``original OWR report'': \url{https://doi.org/10.4171/owr/2016/60}.
\item Literature assessment checked at \texttt{2026-08-21}: the supplied dated curation labels the problem ``partially solved'' and reports that every real closed field has an integer part, that models of true arithmetic occur as integer parts of suitable real exponential fields, and that the specific surreal-number and normal-integer-part questions remain unresolved. These are supplied literature assessments, not independently verified mathematical or current-status claims. The background explicitly requests independent verification of time-sensitive status and cited claims.
\item The supplied registry labels the statement \texttt{corrected\_verified} and reports that the official report passage was checked and rewritten as a self-contained problem preserving its mathematical content. This is a supplied verification claim; no source was visited or independently checked for this draft.
\end{itemize}
\end{document}
